\section{Formal Definition of Concept Activation Regions}
\label{app:concept_regions}
This appendix provides the formal construction of the concept activation function $\boldsymbol{\Phi}_\theta$ and activation region $\mathcal{R}_c$ used throughout~\cref{sec:tradeoff}. The compressed version in the main text specifies these informally; the definition here is the formal object used in the proofs of~\cref{thm:tradeoff}.

\begin{definition}[Concept activation function and region]
\label{def:activation}
Fix a UNet layer $\ell^*$ and denoising timestep $t^* \in \{1,\ldots,T\}$. Let $p \in \mathcal{P}$ denote a text prompt and $\mathbf{e}_p = f_\psi(p)$ be its text embedding. The \emph{concept activation function} of the pretrained model is
\begin{equation}
    \boldsymbol{\Phi}_{\theta, \ell^*,t^*}(p) \;=\; \mathrm{A}_{\theta}(\mathbf{e}_p,\ell^*,t^*) \in \mathcal{H},
\end{equation}
where $\mathrm{A}_{\theta}(\mathbf{e}_p, \ell^*, t^*)$ denotes the cross-attention output 
%\varun{perhaps avoid cal A? which is commonly used for algorithm or adversary? can use A?}
of $\epsilon_\theta$ at layer $\ell^*$ and timestep $t^*$, and $\mathcal{H} := \mathbb{R}^{d_{\ell^*}}$ is the corresponding activation space. 
Throughout, we suppress the dependence on $\ell^*$ 
and $t^*$ and write $\boldsymbol{\Phi}_\theta(p)$ for brevity.
Analogously, $\boldsymbol{\Phi}_{\hat{\theta}}(p)$ denotes the activation function of the CEM $\mathcal{D}_{\hat\theta}$.

For a concept $c \in \mathcal{C}$, its \emph{activation region} is
\begin{equation}
    \mathcal{R}_c
    \;=\;
    \Bigl\{
        \boldsymbol{\Phi}_{\theta}(p)
        \;\Big|\;
        p \in \mathcal{P},\;
        \mathbb{O}_{\mathcal{X}}\!\left(\mathbf{x}_{\theta}(p),\, c\right) = 1
    \Bigr\}
    \;\subseteq\; \mathcal{H},
\end{equation}
% \ali{this also depends on $\ell^*,t^*$} 
where $\mathbb{O}_{\mathcal{X}} : \mathcal{X} \times \mathcal{C} \to \{0,1\}$ is an oracle indicating whether concept $c$ is present in image $\mathbf{x}$, and $\mathbf{x}_{\theta}(p) \sim p_{\theta}(\cdot \mid p)$ denotes an image sampled from $\mathcal{D}_\theta$ conditioned on prompt $p$.
\end{definition}

\paragraph{On the oracle classifier.}
The predicate $c \in \text{image}(\mathbf{x})$ is specified as an oracle classifier rather than a specific model because our results are independent of the particular choice. Any reasonable concept detector induces a concept region $\mathcal{R}_c$, and the tradeoff in~\cref{thm:tradeoff} holds with respect to whichever classifier is fixed. In~\cref{sec:experiment}, we instantiate this oracle with a CLIP-based classifier, detailed in~\cref{app:evaluation_setup}. 